%% ch5_implementation.tex  —  Chapter 5: Implementation
%% HSP-Agile Thesis, East China Normal University, 2026
%% Ye Xujun

\chapter{Implementation}
\label{ch:implementation}

This chapter is engineering, not re-derivation.
Reviewers who already accept the method need to know which module owns which
failure mode, what can be swapped for ablation, and where latency is spent.
Algorithms stay in Chapter~\ref{ch:method}; here the concern is modules,
control flow, and trade-offs.
On \code{classify\_signal}, the same modules emit the witness $(-3,-10)$,
serialise the Semantic Feedback IR record, and decide $\mathrm{Accept}$.

Notation: $t = (\mathcal{S}_t, \mathcal{X}_t, \mathcal{Y}_t)$, candidate~$c$,
$\mathrm{Conf}(c,t)$, $\mathrm{Accept}(c,t)$ from~\eqref{eq:accept}.

% ---------------------------------------------------------------------------
\section{System Architecture Overview}
\label{sec:impl:architecture}
% ---------------------------------------------------------------------------

Why separate modules?
Because Stages~1--4 must be ablatable and auditable without rewriting the
rest (Chapter~\ref{ch:experiments}).
Ingestion, orchestration, and external services therefore stay apart.

For the monitoring alert \code{classify\_signal}, the path is: adapter
$\to$ SpecIR $\to$ LLM candidate $\to$ Z3 witness $(-3,-10)$ $\to$ Semantic
Feedback IR $\to$ pattern screen $\to$ $\mathrm{Accept}$.
Figure~\ref{fig:component-arch} is that layout.
E8 adapters (Mini-Z, Mini-StateMachine) feed the same SpecIR path; typed
result objects keep stages isolatable.

\begin{figure}[t]
  \centering
  \IfFileExists{diagrams/puml/rendered/component_arch.pdf}{%
    \pumlfig{diagrams/puml/rendered/component_arch}%
  }{%
    \fbox{\parbox{0.90\linewidth}{\centering\vspace{1.2cm}
      \textit{[Figure pending render: component\_arch.pdf]}\\
      HSP-Agile system component architecture.
      \vspace{1.2cm}}}%
  }
  \caption{\textbf{Key observation:} Stages~1--4 are isolatable modules under one
    orchestrator, so ablations swap a stage without rewriting the rest.
    \textbf{How to read:} pipeline core (centre) vs.\ external LLM/SMT (below)
    vs.\ experiment orchestration (above); software view of
    Figure~\ref{fig:sgdp-overview}.}
  \label{fig:component-arch}
\end{figure}

\texttt{PipelineRunner} owns the loop: immutable \texttt{PipelineConfig},
delegation to \texttt{ECNUClient}, \texttt{FormalChecker}, and
\texttt{PatternMatcher} (Figure~\ref{fig:system-class}).
Value objects carry serialisable audit traces.

\begin{figure}[t]
  \centering
  \IfFileExists{diagrams/puml/rendered/system_class.pdf}{%
    \pumlfig{diagrams/puml/rendered/system_class}%
  }{%
    \fbox{\parbox{0.90\linewidth}{\centering\vspace{1.2cm}
      \textit{[Figure pending render: system\_class.pdf]}\\
      HSP-Agile system class diagram.
      \vspace{1.2cm}}}%
  }
  \caption{\textbf{Key observation:} one \texttt{run(task)} entry point composes
    Stages~1, 2, and~4 behind Algorithm~\ref{alg:sgdp}.
    \textbf{How to read:} \texttt{PipelineRunner} at the hub; clients and
    checkers as leaves.}
  \label{fig:system-class}
\end{figure}

\paragraph{What each subsystem prevents.}
Adapters keep notation changes out of the core: SOFL, Mini-Z, and
Mini-StateMachine each implement \texttt{SpecAdapter.parse() $\to$ SpecIR},
then lower via \texttt{FSFLowerer}.
The pipeline runner realises Algorithm~\ref{alg:sgdp}, including early exit on
$\mathrm{Accept}(c_k,t)=1$ and best-effort fallback when $k > K$.
The formal checker builds Z3 queries, executes candidates in a sandbox, and
assembles counterexamples.
The pattern guard applies catalogue $\mathcal{P}$ (Table~\ref{tab:patterns})
with AST and regex detectors.
The experiment orchestrator schedules mode--task pairs across a process pool
with resume-safe logs.

% ---------------------------------------------------------------------------
\section{Specification Ingestion Layer}
\label{sec:impl:ingestion}
% ---------------------------------------------------------------------------

The checker must never see raw SOFL, Mini-Z, or Mini-StateMachine text.
Ingestion therefore parses each surface notation into SpecIR
(Chapter~\ref{ch:formalization}) and lowers it to the FSF-shaped
\texttt{TaskSpec} consumed by Stages~2--4.
Three phases do the work: notation-specific parsing via a registered
\texttt{SpecAdapter}; canonicalisation into ordered \texttt{GuardedCase}
records; and lowering via \texttt{FSFLowerer}.
The primary SOFL path uses the Agile-SOFL
toolchain~\citep{li2022agilesofl,liu2014sofl}; E8 adapters bypass SOFL
syntax while producing the same SpecIR structure.
On the SOFL path, the parser classes in Figure~\ref{fig:fsf-parser-class}
are what keep surface syntax from leaking into the checker.

\subsection{SpecIR Adapter Registry}
\label{sec:impl:specir-adapters}

The adapter registry (\path{src/adapters/__init__.py}) maps notation names to
concrete \texttt{SpecAdapter} subclasses.
Each adapter implements \texttt{parse(spec\_text, task\_id) $\to$ SpecIR},
populating \texttt{GuardedCase} objects with normalised \texttt{GuardAtom}
predicates, postconditions, and surface-text fields for prompting.
\texttt{SpecIR.to\_dict()} and \texttt{SpecIR.from\_dict()} round-trip through
\path{schemas/spec_ir.schema.json}, enabling audit logging and cross-notation
benchmark interchange.

\texttt{FSFLowerer.lower(spec)} projects \texttt{SpecIR} onto the legacy
\texttt{TaskSpec} schema: ordered scenario list, I/O signatures, prompt text,
and metadata required by the pattern guard.
\texttt{PipelineRunner.run()} accepts either a \texttt{TaskSpec} dict or a
\texttt{SpecIR} instance (\path{src/pipeline/task_input.py}), normalising
through the lowerer when needed.
This separation keeps oracle construction, witness generation, and repair
feedback notation-agnostic while preserving backward compatibility with the
120-task FSF benchmark and existing experiment scripts.

\begin{figure}[t]
  \centering
  \IfFileExists{diagrams/puml/rendered/fsf_parser_class.pdf}{%
    \pumlfig{diagrams/puml/rendered/fsf_parser_class}%
  }{%
    \fbox{\parbox{0.90\linewidth}{\centering\vspace{1.2cm}
      \textit{[Figure pending render: fsf\_parser\_class.pdf]}\\
      FSF specification parsing and test-case generation class diagram.
      \vspace{1.2cm}}}%
  }
  \caption{FSF parsing and downstream translation (Stage~2 inputs of
           Chapter~\ref{ch:method}).
           For \code{classify\_signal}, the bridge yields a three-scenario
           \texttt{TaskSpec}; the Z3 translator emits \texttt{ConcreteCase}
           witnesses such as $(-3,-10)$.}
  \label{fig:fsf-parser-class}
\end{figure}

\subsection{Parsing and TaskSpec Construction}

The Agile-SOFL parser performs lexical and syntactic analysis of module
definitions, extracting process and function bodies together with their
associated FSF blocks.  For each process~$p$ that declares a non-empty FSF
specification, the bridge constructs a \texttt{TaskSpec} record containing:
\begin{itemize}
  \item \textbf{Task identifier} --- a stable string key derived from module
        and process names;
  \item \textbf{Function signature} --- ordered input parameter names and
        types, output variable declarations, and the callable name used in
        candidate execution;
  \item \textbf{Scenario list} --- an ordered sequence of
        \texttt{FSFScenario} objects, one per FSF clause;
  \item \textbf{External variable map} --- read/write declarations from the
        \texttt{ext rd}/\texttt{ext wr} blocks of the process body;
  \item \textbf{Prompt specification} --- a human-readable rendering of the
        FSF text, injected verbatim into the LLM synthesis prompt.
\end{itemize}

Each \texttt{FSFScenario} stores three fields aligned with
Definition~\ref{def:scenario}: a guard condition string (\texttt{test}),
an output assignment block (\texttt{def}), and a scenario kind flag
(\texttt{scenario} or \texttt{others}).  Scenario indices preserve the
first-match ordering $s_1 \succ s_2 \succ \cdots \succ s_n$ that governs
the specification oracle~$g_t$, following guarded-command
semantics~\citep{dijkstra1975guarded}.

\subsection{FSF Scenario Extraction}

Scenario extraction walks the parser's concrete syntax tree for FSF
expressions.  A standard FSF block chains scenarios as
\[
  \begin{aligned}
    &\mathit{guard}_1 \;\&\&\; \mathit{assign}_1
    \;\|\|\; \mathit{guard}_2 \;\&\&\; \mathit{assign}_2
    \;\|\|\; \cdots \\
    &\|\|\; \mathit{others} \;\&\&\; \mathit{default\_assignments}.
  \end{aligned}
\]
The extractor splits on \code{||} to obtain individual scenarios, then
separates each scenario's guard predicate from its output mapping at the
\code{\&\&} delimiter.

Guards are stored in a normalised relational atom syntax (e.g.\
\texttt{a gt 10} for $a > 10$) that is shared across three consumers: the
LLM prompt renderer, the Z3 translator, and the pattern guard's
guard-comparison detectors.  Output assignments are parsed into dictionaries
mapping output variable names to integer constants or symbolic references to
input variables, enabling both oracle evaluation and prompt signature
injection.

\subsection{External Variable Handling}

Agile-SOFL processes may declare \emph{external variables}---module-level
state that a process reads or writes during execution.  The ingestion layer
records these as a typed map $\mathit{ext} : \mathit{Name} \to
\{\textsc{rd}, \textsc{wr}\}$.  When the prompt is assembled, readable
external variables are listed alongside input parameters so that the LLM is
aware of the full variable environment.  The pattern guard subsequently
verifies that declared external variables appear in the candidate's name
reference set (pattern RF10), closing a common gap where the model ignores
SOFL state variables entirely.

\subsection{Output Inference}

Output variable names are inferred from two sources: the explicit output
parameter declarations in the process signature, and the left-hand sides of
assignment atoms in each scenario's \texttt{def} block.  The union of these
sets defines the output domain~$\mathcal{Y}_t$ for conformance comparison.
Requiring the candidate to return a dictionary keyed by exactly these names
eliminates ambiguity during field-by-field oracle comparison in Stage~2 and
enables RF04 (missing output field) detection in Stage~4.

% ---------------------------------------------------------------------------
\section{LLM Integration Layer}
\label{sec:impl:llm}
% ---------------------------------------------------------------------------

Synthesis is the only non-deterministic stage; everything else must be
replayable.
The LLM layer wraps an OpenAI-compatible gateway
(\texttt{ECNUClient}, \path{src/llm/ecnu_client.py}) for auth, retry, timeout,
token accounting, and audit logs---the name is historical, not a coupling.

\subsection{LLM Client Adapter Design}

The LLM client adapter (\texttt{ECNUClient}) exposes a single \texttt{chat(messages)}
method that accepts a list of role--content message pairs and returns the
model's completion text together with usage metadata (prompt tokens,
completion tokens, wall-clock latency).  Three resilience mechanisms are
built in:

\begin{enumerate}
  \item \textbf{Exponential backoff retry.}  Transient network failures and
        HTTP 5xx responses trigger up to three retry attempts with backoff
        interval $2^{\mathit{attempt}}$ seconds, preventing a single timeout
        from aborting an entire experiment batch.
  \item \textbf{Request timeout.}  Each API call is bounded by a
        90-second wall-clock timeout.  Calls exceeding this limit are
        treated as failures and retried; persistent timeout triggers
        escalation to the repair loop's best-effort exit.
  \item \textbf{Model selection.}  All production experiments use
        Qwen2.5-27B-Instruct (\textsf{Qwen-27B}), which provides strong code
        generation capability on structured specification prompts while
        remaining accessible through the standard chat-completions API
        surface, consistent with GPT-class code models~\citep{openai2023gpt4,chen2021codex}.
\end{enumerate}

Every successful call is optionally serialised to a JSON audit log containing
the full message history, response text, token counts, and task metadata
(task identifier, mode, attempt index).  This design supports cost tracking
and enables manual inspection of failure cases without re-invoking the API.

\subsection{Prompt Template Structure}

Prompt construction follows the Stage~1 template of
Section~\ref{subsec:method:prompt}.  The adapter does not embed domain logic; instead,
the pipeline runner assembles a two-message chat completion request:

\begin{itemize}
  \item \textbf{System message.}  Establishes the role (expert Agile-SOFL
        Python programmer), output format constraints (single fenced code
        block, dictionary return type, integer-only arithmetic), and the
        critical instruction to evaluate FSF guards in listed order.
  \item \textbf{User message.}  Concatenates the rendered FSF specification
        (\texttt{promptSpec}), the extracted function signature, the list of
        required output keys, and---on repair iterations $k > 1$---the
        verbalised feedback from the previous attempt.
\end{itemize}

System messages stay fixed; user messages grow with counterexamples
\cite{chen2021codex,lynn2024verifierloop}.

\subsection{Temperature and Top-$p$ Configuration}

$T_\text{gen}{=}0.7$ for first attempts, $T_\text{repair}{=}0.0$ for repair,
$\mathit{top\_p}{=}0.95$---explore once, then exploit under a deterministic
checker (Chapter~\ref{ch:experiments}).

\subsection{Semantic Feedback IR and Two-Layer Repair Prompting}
\label{sec:impl:repair-feedback}

When iteration $k$ fails formal checking or pattern screening, the repair engine
constructs a \texttt{SemanticFeedbackIR} (\path{src/repair/feedback_ir.py})
from checker counterexamples and the active \texttt{TaskSpec}.
Each record serialises to JSON per
\texttt{schemas/semantic\_feedback.schema.json} \emph{before} any natural-language
rendering, establishing a machine-readable intermediate layer distinct from the
repair prompt:
\begin{enumerate}
  \item \textbf{IR construction} --- nine typed fields (\texttt{violation\_type},
        \texttt{scenario\_index}, \texttt{constraint\_text}, \texttt{inputs},
        \texttt{expected}, \texttt{observed}, \texttt{reason}, \texttt{priority},
        \texttt{suggested\_fix}) populated from witness and scenario metadata;
  \item \textbf{JSON serialisation} --- \texttt{SemanticFeedbackIR.to\_json\_list()}
        emits schema-compliant records logged in \texttt{PipelineResult.semantic\_feedback};
  \item \textbf{Prompt projection} --- \texttt{FeedbackRenderer.render()} maps the
        same IR to one of three surfaces: \texttt{test\_only} (mode~B2 / E6-A),
        \texttt{test\_expected} (E6-B), or \texttt{semantic\_ir} (mode~M / E6-C).
\end{enumerate}

Pattern violations are appended as a semicolon-separated list of
\texttt{pattern\_id:name} pairs (capped at five entries).
The complete feedback block is injected into the user message under a
\emph{``Previous attempt failed''} header.
This IR$\to$JSON$\to$Prompt pipeline proves that Semantic Feedback is a genuine
intermediate representation, not merely a prompt template; E6 varies only the
renderer while holding IR construction fixed.

\subsection{Output Parsing}

The model response is scanned for fenced Python code blocks.  When multiple
blocks are present, the last block is taken as the candidate implementation
(a conservative choice that discards explanatory preamble).  If no fenced
block is found, a heuristic line filter retains lines beginning with
\texttt{def}, \texttt{import}, \texttt{return}, or continuation lines of an
already-started block.  The extracted source is compiled via Python's
\texttt{compile()} primitive; compilation failure yields a sentinel
candidate~$\bot$ that fails all conformance cases with score zero, triggering
repair on the next iteration without crashing the pipeline.

% ---------------------------------------------------------------------------
\section{Formal Conformance Checker}
\label{sec:impl:checker}
% ---------------------------------------------------------------------------

Acceptance must not rest on lucky unit tests.
The formal checker is the deterministic oracle that grounds release decisions
in SMT-backed witnesses rather than heuristic
sampling~\citep{claessen2000quickcheck}.
Given candidate~$c$ and task~$t$, it computes $\mathrm{Conf}(c,t)$ and
assembles the counterexample set $\mathit{cexs}$ for repair.
Case generation and evaluation are the two partitions in
Figure~\ref{fig:checker-activity}.

\begin{figure}[t]
  \centering
  \IfFileExists{diagrams/puml/rendered/checker_activity.pdf}{%
    \pumlfig{diagrams/puml/rendered/checker_activity}%
  }{%
    \fbox{\parbox{0.90\linewidth}{\centering\vspace{1.2cm}
      \textit{[Figure pending render: checker\_activity.pdf]}\\
      Formal conformance checker activity diagram.
      \vspace{1.2cm}}}%
  }
  \caption{Activity diagram of the formal conformance checker (Stage~2,
           Algorithm~\ref{alg:checker}).
           Case generation and candidate evaluation are separate partitions;
           counterexamples feed Stage~3 repair on the next synthesis iteration.}
  \label{fig:checker-activity}
\end{figure}

\subsection{FSF-to-Z3 Translation Strategy}

Guard predicates from each \texttt{FSFScenario} are translated into Z3
integer constraints following the priority encoding of
equation~\eqref{eq:priority}.  The translator
maintains a symbol table mapping each variable name in $\mathcal{X}_t \cup
\mathcal{Y}_t$ to a Z3 \texttt{Int} constant.  Relational atoms in the
normalised syntax (\texttt{gt}, \texttt{lt}, \texttt{eq}, etc.) are mapped
to the corresponding Z3 arithmetic predicates.  Conjunctions and disjunctions
in guard expressions are recursively translated to \texttt{And} and
\texttt{Or} nodes.

To prevent unbounded solver search, integer variables are constrained to a
finite domain $[-5, 20]$ during model generation.  This domain is sufficient
for the benchmark tasks in Chapter~\ref{ch:experiments}, whose inputs lie in
$[0, 100]^5$ but whose guard thresholds are small integers; the constraint
primarily serves as a solver termination guarantee rather than a semantic
restriction on $\mathcal{X}_t$.

\subsection{Scenario Precedence Encoding}

First-match semantics require that a test case for scenario~$s_i$ activate
$s_i$ and no higher-priority scenario.  The implementation encodes this as
the priority-constrained formula $\Phi_i$ from
equation~\eqref{eq:priority}:
\[
  \Phi_i \;\equiv\; \widehat{g}_i \;\land\;
    \bigwedge_{j=1}^{i-1} \lnot\,\widehat{g}_j.
\]
For the final \texttt{others} scenario, $\widehat{g}_n$ is encoded as the
negation of the disjunction of all preceding guards, matching
Definition~\ref{def:oracle}.  Each $\Phi_i$ is submitted to an independent
Z3 \texttt{Solver} instance, ensuring that solver state from one scenario
does not contaminate another.

\subsection{Concrete Case Generation Algorithm}

Algorithm~\ref{alg:case-gen} presents the case generation procedure used
both during repair-loop iterations and for final strict evaluation.  The
algorithm supports a \texttt{max\_cases} parameter that controls how many
witnesses are produced per scenario through model blocking: after extracting
a satisfying model, the solver adds a disjunctive constraint excluding that
model, then re-queries for an alternative witness.

\begin{algorithm}[t]
\caption{Z3-Based Concrete Case Generation}
\label{alg:case-gen}
\begin{algorithmic}[1]
\Require ordered scenarios $\mathcal{S}_t = \langle s_1,\ldots,s_n\rangle$,
         variable set $V$, budget $N$
\Ensure case set $\mathcal{D}(t)$ with $|\mathcal{D}(t)| \le N$
\State $\mathcal{D}(t) \leftarrow \emptyset$
\State $\mathit{sym} \leftarrow \{v \mapsto \ProcCall{Z3Int}{v} : v \in V\}$
\For{$i = 1$ \textbf{to} $n$}
  \State $\Phi_i \leftarrow \ProcCall{GuardToZ3}{g_i} \land
         \bigwedge_{j<i} \lnot\,\ProcCall{GuardToZ3}{g_j}$
         \Comment{eq.~\eqref{eq:priority}}
  \State $\mathit{solver} \leftarrow \ProcCall{FreshSolver}{}$
  \State $\ProcCall{Add}{\mathit{solver},\,\Phi_i}$
  \State $\ProcCall{Add}{\mathit{solver},\,
         \bigwedge_{v \in V} -5 \le \mathit{sym}[v] \le 20}$
  \State $\mathit{attempts} \leftarrow 0$
  \While{$\ProcCall{Check}{\mathit{solver}} = \mathit{sat}$
         \textbf{and} $\mathit{attempts} < 3$
         \textbf{and} $|\mathcal{D}(t)| < N$}
    \State $x \leftarrow \ProcCall{ExtractModel}{\mathit{solver}}$
    \State $y \leftarrow \phi_i(x)$ \Comment{evaluate output mapping}
    \State $\mathcal{D}(t) \leftarrow \mathcal{D}(t) \cup \{(x, y, i)\}$
    \State $\ProcCall{Add}{\mathit{solver},\,
           \bigvee_{v \in V} \mathit{sym}[v] \ne x[v]}$
    \State $\mathit{attempts} \leftarrow \mathit{attempts} + 1$
  \EndWhile
\EndFor
\State \Return $\mathcal{D}(t)$
\end{algorithmic}
\end{algorithm}

\subsection{Candidate Execution and Conformance Scoring}

After case generation, the checker compiles the candidate into an isolated
namespace, invokes it with each witness input as keyword arguments, and
compares the returned dictionary against the oracle output field-by-field
using integer equality.  Runtime exceptions are recorded as counterexamples
with an empty actual-output field and an error annotation.

The conformance score follows equation~\eqref{eq:conformance}:
$\mathrm{Conf}(c,t) = |\mathit{passed}| / |\mathcal{D}(t)|$.
A candidate achieves \emph{strict formal pass} when
$\mathrm{Conf}(c,t) = 1$ and $|\mathcal{D}(t)| > 0$.

\subsection{Strict vs.\ Standard Evaluation Case Counts}

The checker is invoked at two distinct granularities, reflecting a deliberate
latency--coverage trade-off:

\begin{itemize}
  \item \textbf{Standard (repair-loop) evaluation} uses a budget of
        16 cases per task during repair iterations.
  \item \textbf{Strict (final) evaluation} uses a budget of
        64 cases per task for the post-pipeline strict conformance score
        reported in Chapter~\ref{ch:experiments}.
\end{itemize}

The 16-case internal budget and 64-case strict budget are configuration
parameters on \texttt{PipelineConfig}; neither is hard-coded, allowing
sensitivity analysis over case-count settings (reported in
Section~\ref{sec:results:sensitivity} of Chapter~\ref{ch:results}).

% ---------------------------------------------------------------------------
\section{Counterexample-Guided Repair Engine}
\label{sec:impl:repair}
% ---------------------------------------------------------------------------

Without careful state management, a repair loop either loops forever or
forgets why the previous candidate failed.
The repair engine adapts CEGIS~\citep{solarlezama2008sketching} to an
LLM synthesiser: it keeps the feedback buffer, increments the attempt
counter, and decides whether to continue or exit with a best-effort
candidate.
Figure~\ref{fig:pipeline-state} makes those control states explicit; the
data that flows on each repair arc is Figure~\ref{fig:repair-feedback}
(Chapter~\ref{ch:method}).

\begin{figure}[t]
  \centering
  \IfFileExists{diagrams/puml/rendered/pipeline_state.pdf}{%
    \pumlfig{diagrams/puml/rendered/pipeline_state}%
  }{%
    \fbox{\parbox{0.90\linewidth}{\centering\vspace{1.2cm}
      \textit{[Figure pending render: pipeline\_state.pdf]}\\
      HSP-Agile pipeline state machine.
      \vspace{1.2cm}}}%
  }
  \caption{Pipeline state machine (implements Stage~3 of
           Chapter~\ref{ch:method}).  The repair state aggregates
           counterexamples and pattern hits into the feedback buffer before
           returning to synthesis.  Acceptance requires both formal
           conformance and a clear high-severity pattern screen.}
  \label{fig:pipeline-state}
\end{figure}

\subsection{Feedback State Machine}

The pipeline state machine (Figure~\ref{fig:pipeline-state}) has six states:
\textsc{Init}, \textsc{Synthesizing}, \textsc{FormalChecking},
\textsc{PatternChecking}, \textsc{Repairing}, and terminal states
\textsc{Accepted} or \textsc{Exhausted}.  Transitions are governed by the
intermediate predicates:
\begin{align*}
  \textsc{FormalChecking} \to \textsc{PatternChecking}
    &\quad\text{when } \mathrm{Conf}(c_k, t) = 1, \\
  \textsc{FormalChecking} \to \textsc{Repairing}
    &\quad\text{when } \mathrm{Conf}(c_k, t) < 1, \\
  \textsc{PatternChecking} \to \textsc{Accepted}
    &\quad\text{when } {\textstyle\sum_{p \in \mathcal{P}^H}} \delta_p(c_k,t) = 0, \\
  \textsc{PatternChecking} \to \textsc{Repairing}
    &\quad\text{when any } p \in \mathcal{P}^H \text{ fires}.
\end{align*}

The \textsc{Repairing} state assembles $\mathit{feedback}_k$ from
$\mathit{cexs}_k$ and $\mathit{hits}_k$, increments $k$, and returns to
\textsc{Synthesizing} if $k \le K$.

\subsection{Budget $K$ and Early Termination}

The repair budget $K$ is a hard upper bound on LLM invocations per task,
enforcing design principle~P2 from Chapter~\ref{ch:method}.  In the
experimental configuration, $K = 3$: the pipeline makes at most three LLM
calls regardless of task difficulty.  Two early-exit conditions can
terminate the loop before $K$ is reached:

\begin{enumerate}
  \item \textbf{Acceptance exit.}  When $\mathrm{Accept}(c_k, t) = 1$,
        the current candidate is returned immediately with status
        \textsc{Accept}.  No further LLM calls are made.
  \item \textbf{Non-repairable exit.}  When \texttt{enable\_repair} is
        \textit{false} (mode A3), the pipeline performs a single synthesis
        call and proceeds directly to evaluation without feedback, regardless
        of the conformance outcome.
\end{enumerate}

If the budget is exhausted without acceptance, the pipeline returns the
candidate with the highest observed $\mathrm{Conf}(c_k, t)$ across all
iterations, annotated with status \textsc{Best-Effort}.  This fallback
ensures that the system always produces an auditable artefact even on tasks
where the LLM cannot fully repair within the allotted budget.

Algorithm~\ref{alg:repair-loop} formalises the repair engine control logic.

\begin{algorithm}[t]
\caption{Counterexample-Guided Repair Loop}
\label{alg:repair-loop}
\begin{algorithmic}[1]
\Require task $t$, config with budget $K$, flags
         $(\mathit{enable\_formal}, \mathit{enable\_patterns}, \mathit{enable\_repair})$
\Ensure accepted candidate $c^*$, status $\in \{\textsc{Accept}, \textsc{Best-Effort}\}$
\State $\mathit{feedback} \leftarrow \emptyset$;\quad
       $\mathit{best} \leftarrow (\bot, 0)$
\For{$k = 1$ \textbf{to} $K$}
  \State $c_k \leftarrow \ProcCall{LLM}{t,\, \mathit{feedback}}$
  \If{$\mathit{enable\_formal}$}
    \State $(\mathit{conf}_k, \mathit{cexs}_k) \leftarrow
           \ProcCall{FormalCheck}{c_k,\, t,\, 16}$
  \Else
    \State $(\mathit{conf}_k, \mathit{cexs}_k) \leftarrow
           \ProcCall{UnitTestScore}{c_k,\, t}$
  \EndIf
  \If{$\mathit{enable\_patterns}$}
    \State $\mathit{hits}_k \leftarrow \ProcCall{PatternGuard}{c_k,\, t}$
  \Else
    \State $\mathit{hits}_k \leftarrow \emptyset$
  \EndIf
  \If{$\mathit{conf}_k > \mathit{best}.\mathit{score}$}
    \State $\mathit{best} \leftarrow (c_k, \mathit{conf}_k)$
  \EndIf
  \If{$\ProcCall{Accept}{c_k,\, t} = 1$}
    \State \Return $c_k$,\enspace \textsc{Accept}
  \EndIf
  \If{$\mathit{enable\_repair}$}
    \State $\mathit{feedback} \leftarrow
           \ProcCall{BuildFeedback}{\mathit{cexs}_k,\, \mathit{hits}_k,\, k}$
           \Comment{Stage 3}
  \EndIf
\EndFor
\State \Return $\mathit{best}.\mathit{candidate}$,\enspace \textsc{Best-Effort}
\end{algorithmic}
\end{algorithm}

% ---------------------------------------------------------------------------
\section{Pattern Guard Module}
\label{sec:impl:patterns}
% ---------------------------------------------------------------------------

A candidate can match every witness and still hard-code success or drop an
output field that the finite case set never exercised.
Stage~4 therefore screens structural anti-patterns before release.
The module evaluates catalogue $\mathcal{P}$ (RF01--RF14) and returns
$\mathrm{Hits}(c,t)$ as in Section~\ref{subsec:method:detection}, using a
hybrid AST/regex strategy in the spirit of FindBugs
\citep{hovemeyer2004findbugs,ayewah2008findbugs}.

\subsection{Pattern Catalogue and High-Severity Subset}

Table~\ref{tab:patterns} (Chapter~\ref{ch:method}) lists the full catalogue.
Patterns are classified by severity:
\emph{critical} (RF07), \emph{high}, \emph{medium}, and \emph{low}.
The high-severity subset $\mathcal{P}^H$ used in
equation~\eqref{eq:accept} comprises nine patterns whose detection
immediately blocks acceptance:
\[
  \mathcal{P}^H = \{\text{RF01}, \text{RF02}, \text{RF04}, \text{RF05},
    \text{RF06}, \text{RF07}, \text{RF09}, \text{RF11}, \text{RF13}\}.
\]
Medium- and low-severity hits (RF03, RF08, RF10, RF12, RF14) are recorded in
the audit trace and included in repair feedback but do not gate acceptance.

\subsection{AST-Based and Regex-Based Matching Strategy}

Detection follows a hybrid strategy keyed by pattern type:

\begin{itemize}
  \item \textbf{AST-based detectors} parse the candidate into a Python
        abstract syntax tree and traverse it for structural properties.
        RF04 counts return-statement dictionary keys against the declared
        output domain.  RF07 detects unconditional success returns on code
        paths with no conditional branches.  RF09 compares the number of
        \texttt{if} branches against the non-default scenario count.
        RF13 identifies empty \texttt{else} bodies containing only
        \texttt{pass}.
  \item \textbf{Regex-based detectors} scan the raw source text for surface
        patterns.  RF03 matches hardcoded numeric literals not present in
        the specification.  RF08 detects float-return annotations where
        integer output is required.  RF14 compares docstring scenario
        references against actual guard conditions.
  \item \textbf{Cross-reference detectors} combine extracted FSF guard
        predicates with the candidate's conditional structure.  RF01, RF02,
        RF05, RF06, and RF11 compare threshold literals and comparison
        operators in the implementation against the specification's guard
        atoms, detecting missing preconditions, wrong default branches,
        swapped comparisons, and off-by-one boundary errors.
\end{itemize}

Each detector is a pure function $\delta_p : \mathcal{C} \times \mathcal{T}
\to \{0,1\}$; the guard module evaluates all 14 detectors sequentially and
collects firing patterns into a list of \texttt{PatternHit} records, each
carrying pattern identifier, severity, source location, and evidentiary
message.

\subsection{Conjunctive Acceptance Gate Role}

The pattern guard is the second conjunct of $\mathrm{Accept}(c,t)$:
\[
  \mathrm{Accept}(c, t) =
    \mathbf{1}\!\left[\mathrm{Conf}(c, t) = 1\right]
    \land
    \mathbf{1}\!\left[\sum_{p \in \mathcal{P}^H} \delta_p(c, t) = 0\right].
\]
In the implementation, the guard is consulted \emph{after} formal checking
on each iteration.  A candidate that passes all Z3-generated cases but
triggers RF07 (unconditional success) is routed to the repair state rather
than accepted, and the pattern violation is verbalised alongside any
counterexamples in the feedback buffer.  This ordering ensures that the
formal checker and pattern guard operate as complementary filters rather
than redundant substitutes, as argued in
Section~\ref{subsec:conformance-acceptance}.

% ---------------------------------------------------------------------------
\section{Experiment Orchestration}
\label{sec:impl:orchestration}
% ---------------------------------------------------------------------------

An 840-run matrix is useless if a mid-batch crash forces a full restart, or
if Z3 races corrupt results.
The orchestrator therefore prioritises parallel throughput, crash-safe
resumption, and live progress visibility---fitness-function concerns for
sprint-cycle formal methods~\citep{harman2012formalagile,beck2001xp}.

\subsection{Parallel Execution Design}

Each mode--task combination is dispatched as an independent job to a
\emph{process pool}.  Process-level isolation is essential because the Z3
solver maintains global state in its native library; concurrent threads
sharing a single Z3 context have been observed to produce race-induced
\texttt{unknown} results.  By assigning one worker process per CPU core,
the orchestrator achieves near-linear speedup on the 840-run benchmark
(7 modes $\times$ 120 tasks) while preserving solver correctness.

Each worker process constructs its own LLM client adapter (\texttt{ECNUClient}) (with a
process-specific audit log directory), \texttt{PipelineConfig} derived from
the mode identifier, and \texttt{ErrorPreventionPipeline} instance.  After
the pipeline returns, the worker performs a separate strict evaluation pass
with a 64-case budget and optionally evaluates mutation
kill rate~\citep{demillo1978mutation,li2023iet}, then serialises the complete
record to the shared results file.

\subsection{Seven Execution Modes}

Table~\ref{tab:impl-modes} summarises the seven evaluation modes and their
component configurations.  Each mode is realised by setting the three
boolean feature flags (\texttt{enable\_formal}, \texttt{enable\_patterns},
\texttt{enable\_repair}) and associated numeric parameters on
\texttt{PipelineConfig}.

\begin{table}[H]
  \caption{Execution mode configuration matrix.  Flags control which
           pipeline stages are active; numeric parameters tune per-stage
           behaviour.  All LLM-based modes use Qwen2.5-27B-Instruct (\textsf{Qwen-27B})
           with $K = 3$ unless noted.
           $N_{\max}$ is the per-scenario formal case cap.}
  \label{tab:impl-modes}
  \centering
  \thesistable{%
    \footnotesize
    \renewcommand{\arraystretch}{1.15}
    \begin{tabular*}{\linewidth}{@{\extracolsep{\fill}}l l c c c r@{}}
      \toprule
      Mode & Description & Formal & Guard & Repair & $N_{\max}$ \\
      \midrule
      B0 & Reference template  & \textit{eval only} & No  & No  & --- \\
      B1 & One-shot LLM        & No  & No  & No  & --- \\
      B2 & Test-feedback loop  & No  & No  & Yes & 8 \\
      M  & Full HSP-Agile      & Yes & Yes & Yes & 16 \\
      A1 & No formal check     & No  & Yes & Yes & 10 \\
      A2 & No pattern guard    & Yes & No  & Yes & 24 \\
      A3 & No repair loop      & Yes & Yes & No  & 24 \\
      \bottomrule
    \end{tabular*}%
  }
  \par\smallskip
  \raggedright\footnotesize All modes apply strict evaluation with a
  64-case budget after pipeline completion.
  B0 uses a deterministic template renderer; no LLM calls are made.
\end{table}

Mode~\textbf{M} activates all three boolean flags and represents the
complete HSP-Agile configuration.  Ablation modes A1--A3 each disable
exactly one component, enabling the contribution decomposition reported in
Chapter~\ref{ch:experiments}.  Baselines B0--B2 anchor the comparison
against non-LLM, non-iterative, and test-only-feedback approaches
respectively.

\subsection{Resume-Safe Execution and Progress Monitoring}

Long-running experiment batches support \emph{incremental persistence}:
each completed job is appended as a single JSONL record to the results file
immediately upon completion, and the file is flushed to disk before the
worker accepts its next assignment.  On restart, the orchestrator reads
existing records, reconstructs the set of completed task keys, and schedules
only pending jobs.  This design allows interrupted runs---whether from API
rate limits, machine reboot, or manual cancellation---to resume without
re-executing finished work.

A companion \texttt{progress.json} file is updated after every job completion
with: completed count, total count, percentage, elapsed time, estimated time
remaining (computed from the observed completion rate), and the message from
the most recently finished job.  External monitoring tools can poll this
file for live dashboard display without parsing the full results log.

% ---------------------------------------------------------------------------
\section{Design Decisions and Trade-offs}
\label{sec:impl:tradeoffs}
% ---------------------------------------------------------------------------

Every implementation choice trades correctness guarantees against latency and
complexity~\citep{woodcock2009formalsurvey}.
Three decisions dominate the rest.

\subsection{Z3 Over Random Testing}

The formal checker uses Z3~\citep{demoura2008z3} rather than random or
property-based testing~\citep{claessen2000quickcheck,ammann2016testing} for
three reasons grounded in the FSF semantics of Chapter~\ref{ch:formalization}.

First, \emph{guard satisfiability}: random sampling over $\mathcal{X}_t =
\mathbb{Z}^5$ has negligible probability of landing inside a guard region
defined by conjunctions of strict inequalities.  Z3 guarantees that a
witness is found whenever one exists, or reports \textsf{unsat} otherwise.

Second, \emph{precedence correctness}: the priority-constrained formula
$\Phi_i$ produces inputs in scenario~$i$'s \emph{exclusive} region.  A
random input satisfying $g_i$ but also $g_j$ for $j < i$ would be attributed
to scenario~$j$ by the first-match oracle, yielding a misleading test case.
Priority encoding eliminates this ambiguity by construction.

Third, \emph{counterexample minimality}: when a candidate fails, Z3 returns
a concrete witness input that the repair engine can present directly to the
LLM.  Random testing produces failure messages without the scenario-index
context that makes CEGIS feedback effective~\citep{lynn2024verifierloop}.

The cost is solver latency: each Z3 query incurs 50--200\,ms of wall-clock
time, and the strict evaluation pass with 64 cases adds approximately
6--8\,s per task.  The 16-case internal budget during repair limits this
overhead to roughly 2\,s per iteration.

\subsection{Conjunctive Acceptance Over Single-Metric Gating}

An alternative design would accept a candidate when \emph{either}
$\mathrm{Conf}(c,t) = 1$ \emph{or} the pattern guard is clear---a
disjunctive predicate that maximises throughput by accepting on the first
passing criterion.  HSP-Agile rejects this design in favour of the
conjunctive $\mathrm{Accept}(c,t)$ of equation~\eqref{eq:accept} because
the two criteria detect qualitatively different fault classes.

Formal conformance certifies pointwise correctness on the generated test
suite; the pattern guard certifies structural soundness (guard coverage,
output completeness, absence of unconditional returns).  Neither criterion
implies the other: a candidate can achieve $\mathrm{Conf}(c,t) = 1$ by
hard-coding correct outputs at the $n$ witness points while violating RF07
on all other inputs, and conversely a candidate can pass all pattern
detectors while evaluating guards in the wrong order and failing on
overlap-region inputs.

The conjunctive gate is therefore \emph{conservative}: it accepts fewer
candidates than either criterion alone would, but each accepted candidate
satisfies both pointwise and structural requirements.  The experimental
results in Chapter~\ref{ch:experiments} confirm that this selectivity
manifests as a lower strict success rate for mode~M compared to B2, coupled
with higher mean conformance on accepted candidates.

\subsection{Latency vs.\ Quality Trade-off}

The two-tier case-count strategy (16 internal, 64 strict) embodies the
central latency--quality trade-off of the implementation.  Increasing the
internal case count improves per-iteration fault detection at the cost of
proportional Z3 overhead; decreasing it speeds the repair loop but risks
accepting candidates that pass sparse witnesses yet fail on unprobed inputs.

The chosen 16/64 split was calibrated on the development set: 16 cases
provide at least two witnesses per scenario on eight-scenario tasks (sufficient
for counterexample-guided repair), while 64 cases provide eight witnesses
per scenario for final scoring (sufficient to detect ordering faults in
guard-overlap regions).  Sensitivity analysis in
Chapter~\ref{ch:experiments} shows diminishing conformance returns beyond
32 internal cases, validating the chosen operating point.

Similarly, the repair budget $K = 3$ balances correction opportunity against
API cost: empirical observation shows that most repairable faults are resolved
within two iterations, while tasks requiring more than three calls rarely
converge even with additional budget.  The bounded budget ensures
predictable per-task cost in sprint-cycle deployments, consistent with
principle~P2.

% ---------------------------------------------------------------------------
\section*{Chapter Summary}
\addcontentsline{toc}{section}{Chapter Summary}
% ---------------------------------------------------------------------------

Modules own failure modes: adapters keep surface syntax out; Z3 owns
witnesses (16 repair / 64 strict); Semantic Feedback IR owns repair state;
RF01--RF14 own the structural gate; the orchestrator owns crash-safe
throughput.
Whether those choices pay off is empirical---Chapter~\ref{ch:experiments}.
